\documentclass[conference]{IEEEtran}
\usepackage{amsmath,amssymb}
\usepackage{cite}
\usepackage{url}
\usepackage{booktabs}
\usepackage{array}
\usepackage{multirow}

\title{CivicDividendOS: A Factor-Origin Fiscal Ledger and Public Automation Dividend Framework for Human--AI--Robot Economies}

\author{\IEEEauthorblockN{Author Name}
\IEEEauthorblockA{Affiliation\\Email}}

\begin{document}
\maketitle

\begin{abstract}
The rapid diffusion of artificial-intelligence agents, autonomous software, and physical robots is beginning to weaken the historical correspondence between productive work, human wages, payroll taxation, and household income. Existing automation-tax research largely asks whether robots should be taxed and at what rate, while recent AI public-finance research focuses on capital taxation, corporate profits, consumption taxes, transfers, or sovereign wealth funds. These approaches are valuable but leave an unresolved accounting and distribution problem: when an urban service, firm, or digital platform is jointly operated by humans, AI agents, robots, data, and conventional capital, public authorities lack a principled mechanism for determining which factor created economic value, who should legally bear the fiscal obligation, and how automation-generated surplus should be distributed between private owners, workers, citizens, and the state.

This paper proposes \textit{CivicDividendOS}, a novel factor-origin fiscal and distribution framework for economies in which human labour coexists with autonomous digital and embodied productive systems. The framework introduces an \textit{Autonomous Economic Activity Passport} for every significant AI-agent or robotic deployment. The passport does not grant legal personality to machines; instead, it creates an auditable economic identity linking machine activity to its human or corporate owner, operator, jurisdiction, energy use, task class, and generated output. A \textit{Factor-Origin Ledger} then attributes value added across human labour, AI-agent services, robotic services, data contributions, and conventional capital using task-level causal contribution estimates rather than simply imputing a fictional wage to a robot.

Based on this ledger, CivicDividendOS computes a dynamic \textit{Automation Contribution Base} and applies a \textit{Social Automation Contribution} to the legally responsible owner or beneficiary. The rate is endogenous: it rises when autonomous production substitutes for taxable labour, concentrates economic rents, or creates measurable social externalities, and falls through augmentation, reskilling, human-capability, small-firm, and productivity credits when automation complements workers or produces broad social benefit. Collected revenue and selected equity contributions are divided through a \textit{Public Automation Wealth Fund} into four channels: a universal Civic Automation Dividend, targeted worker-transition insurance, a Human Capability Endowment for education and future expertise, and financing for public services. In parallel, a \textit{Human Earned-Income Shield} reduces payroll or labour-income burdens as the taxable labour share contracts.

The framework therefore changes the fiscal unit of analysis from ``taxing robots'' to \textit{taxing the economically attributable automation surplus at the owner/beneficiary level while socializing part of its long-run capital return}. A policy digital twin continuously stress-tests tax incidence, innovation effects, employment, inequality, revenue stability, and cross-border leakage before rates are adjusted. This provides a computational foundation for a tax-and-dividend system that can preserve public revenue and household purchasing power without treating AI agents or robots as fictional legal persons.
\end{abstract}

\begin{IEEEkeywords}
AI taxation, robot tax, automation, fiscal policy, public finance, multi-agent systems,
income distribution, universal basic capital, sovereign wealth fund, future of work,
agentic economy, social dividend.
\end{IEEEkeywords}

\section{Introduction}

Modern tax systems were largely designed around economies in which human work
generates wages, employers remit payroll contributions, households pay income
and consumption taxes, and owners receive taxable profits from capital.
Automation has always complicated this structure, but agentic AI and increasingly
autonomous robots introduce a qualitatively broader challenge. Software agents
can now execute cognitive, administrative, creative, analytical, and commercial
tasks, while robots increasingly perform physical activities previously requiring
human workers.

The economic output of a future firm or city service may therefore arise from a
joint production process involving humans, AI agents, robots, training data,
cloud infrastructure, intellectual property, and conventional capital. If
labour's share of production falls, governments may face a structural mismatch:
the tax base remains strongly linked to wages and payroll while economic value
increasingly accrues to owners of automated productive capital.

The policy question is not simply whether a robot should ``pay tax.'' A robot or
AI agent does not currently possess tax residence, citizenship, consumption
needs, pension rights, or legal responsibility analogous to a human worker.
Current legal analysis generally treats AI and robots as assets, tools, services,
or components of businesses rather than independent taxable persons
\cite{DimitropoulouAIIncome}. The more important problem is how to attribute
economic value produced through autonomous systems and how to distribute a
reasonable share of the resulting surplus without destroying socially valuable
innovation.

This paper proposes CivicDividendOS, a computational framework that combines
factor-origin accounting, adaptive automation contributions, public capital
ownership, citizen dividends, worker transition support, and labour-tax relief.

\section{Literature Review}

\subsection{Automation, Tasks, Wages, and Income Distribution}

A central insight in the economics of automation is that technology changes the
allocation of tasks between labour and capital. Acemoglu and Restrepo show that
automation produces a displacement effect by allowing capital to perform tasks
previously undertaken by labour, while the creation of new tasks can reinstate
labour demand \cite{AcemogluRestrepo2019}. Their task-based perspective provides
an important foundation for analysing fiscal policy because automation affects
not only aggregate productivity but also who receives the associated income.

Empirical research further shows that the labour-market effects of robot
adoption are uneven. Acemoglu and Restrepo document adverse employment and wage
effects in U.S. local labour markets exposed to industrial robots
\cite{AcemogluRestrepo2020}. By contrast, Dauth et al. find that manufacturing
job displacement in Germany was offset by service-sector employment, although
the burden fell disproportionately on younger workers entering the labour market
\cite{Dauth2021}. These findings suggest that automation can generate both gains
and adjustment costs, and that distributional consequences differ across time,
workers, and institutional environments.

Recent AI research broadens these concerns because AI can affect non-routine
cognitive tasks in addition to routine physical or clerical work. The IMF has
emphasized that AI may increase productivity while also widening income and
wealth inequality if gains disproportionately accrue to high-income workers and
capital owners \cite{IMF2024AI}. OECD evidence similarly indicates heterogeneous
effects across occupations and skill groups \cite{OECD2025Korea}.

\subsection{Robot Taxes and Optimal Taxation}

The idea of taxing robots has generated substantial economic and legal debate.
Guerreiro, Rebelo, and Teles develop a quantitative model in which falling
automation costs can sharply increase inequality. They find that robot taxation
can be optimal during the transition while workers exposed to automation remain
in the labour force, but that the optimal tax can disappear after adjustment
\cite{Guerreiro2022}.

Thuemmel studies the joint taxation of robots, other capital, and labour
income. The analysis shows that the optimal treatment of robots depends on
their cost and general-equilibrium effects: subsidies may be desirable when
robots are expensive, while a positive robot tax can emerge as robots become
cheaper and inequality increases \cite{Thuemmel2023}. This demonstrates that
the optimal policy is dynamic rather than a permanent fixed levy.

Acemoglu et al. analyse optimal policy in a task framework and show that reducing
labour taxes while combining capital-tax reductions with targeted automation
taxes can improve employment outcomes relative to uniform capital taxation
\cite{AcemogluOptimalPolicy}. Related work on inefficient automation finds that
firms may automate too rapidly because they do not internalize transition costs
borne by workers, creating an efficiency rationale for slowing automation even
without an explicit redistributive objective \cite{InefficientAutomation}.

Merola reviews taxation as a mechanism for sharing gains from automation and AI,
discussing robot taxes, capital taxation, and redistribution toward workers
\cite{Merola2022}. Zhang similarly analyses automation, wage inequality, and
robot taxation \cite{Zhang2019}.

The literature therefore establishes that some form of differentiated taxation
of automation may be justified under particular conditions. However, a persistent
implementation challenge remains: what exactly constitutes a robot, how should
its contribution be measured, and how can the tax avoid discouraging beneficial
capital deepening?

\subsection{Automation and the Tax Base}

Automation can alter not only employment but also the structure of government
revenue. H{\"o}tte, Theodorakopoulos, and Koutroumpis examine nineteen European
countries and find that the relationship between automation and tax receipts
varies across technologies and phases of diffusion \cite{Hotte2024}. Their
results caution against the simplistic assumption that automation necessarily
causes a permanent decline in total revenue, but they also demonstrate that
automation can change the balance between labour, capital, and consumption tax
bases.

The IMF argues that fiscal policy will need to adapt if AI shifts income toward
capital and away from labour, including stronger social protection and careful
consideration of capital-income taxation \cite{IMF2024Fiscal}. The IMF's 2026
work on the global economic implications of AI similarly emphasizes protecting
tax capacity and fairness if labour income becomes less central
\cite{IMF2026GlobalAI}.

Korinek and co-authors develop a broader public-finance framework for advanced
AI economies, analyzing optimal taxation under scenarios in which capital can
become increasingly dominant \cite{Korinek2026}. This line of work is especially
important because it moves beyond a narrow robot levy toward the redesign of
the fiscal system itself.

\subsection{Taxing AI Rather than Only Robots}

The expansion of generative and agentic AI has renewed interest in taxation
instruments that are not tied to physical robots. Faivre and Cen survey the
viability of AI taxation and distinguish multiple objectives, including
correcting externalities, redistributing uneven gains and costs, and financing
regulatory capacity \cite{FaivreCen2026}. They discuss corporate income and
rent-based taxation, consumption taxes on AI services, and excise-style
instruments tied to particular AI activities.

This work highlights an important design principle: there is unlikely to be one
universal ``AI tax.'' Different harms and policy objectives require different
tax bases. Measurement, tax incidence, international leakage, and innovation
effects are major constraints.

Legal scholarship similarly identifies substantial difficulties in treating
robots or AI systems as independent taxpayers. Current income-tax systems
generally attribute income to owners, firms, or service providers rather than
to the autonomous system itself \cite{DimitropoulouAIIncome}. This creates a
strong case for separating \textit{machine-level economic measurement} from
\textit{legal tax liability}.

\subsection{Universal Basic Income, Social Protection, and Transfers}

One proposed response to automation is redistribution through universal basic
income (UBI) or targeted transfers. Merola notes that revenues from automation
taxes could be used for universal or worker-focused transfers
\cite{Merola2022}. The ILO has examined UBI proposals in relation to social
protection standards and has highlighted financing, adequacy, labour-market,
and institutional questions \cite{ILOUBI}.

Transfers can stabilize household income, but they do not necessarily change
the underlying ownership of productive capital. If AI-driven production becomes
capital intensive, a purely transfer-based model may require governments to
continually tax a concentrated ownership base.

\subsection{Universal Basic Capital and Public Wealth Funds}

An alternative is broader ownership of productive assets. Universal Basic
Capital (UBC) proposes granting citizens financial assets or capital rather than
only periodic income transfers \cite{LeGrand2020}. Progressive sovereign wealth
fund proposals similarly argue that public ownership of capital can generate a
more egalitarian distribution of investment income \cite{Corneo2022}.

The Alaska Permanent Fund provides a long-running real-world example of
converting natural-resource rents into a permanent investment fund whose
earnings support resident dividends \cite{APF}. Although AI is not a natural
resource, the institutional analogy is relevant: temporary or concentrated
economic rents can be partially transformed into broadly owned financial
capital.

Recent public-finance discussions increasingly consider sovereign wealth funds
or citizen-share funds as complements to AI-era tax reform. Brookings argues
that equity-based mechanisms can capture returns that scale with AI's economic
importance \cite{BrookingsTax2026,BrookingsFairness2026}.

\subsection{Data Dividends}

Another literature considers whether individuals should receive dividends from
data used to create AI value. Vincent et al. map design choices for data
dividends and show that apparently reasonable allocation rules can generate
concentrated or demographically unequal outcomes \cite{Vincent2019}. Bax
investigates Shapley-value and Owen-value approaches for estimating the
contribution of individual data samples \cite{Bax2019}.

This literature is useful for value attribution, but data dividends address only
one input into AI production. They do not provide a comprehensive fiscal system
for joint human--AI--robot production.

\subsection{Adaptive and Agent-Based Fiscal Policy}

AI itself is also being investigated as an instrument for fiscal design.
TaxAgent integrates large language models with agent-based economic simulation
to adapt tax policies under heterogeneous household behaviour
\cite{TaxAgent2025}. Other research explores AI-supported detection of tax
loopholes and policy inconsistencies \cite{Fratric2025}. These developments
suggest that future fiscal systems may be increasingly adaptive and
simulation-driven.

However, these systems typically optimize conventional tax structures. They do
not establish an accounting layer that identifies the factor origin of value in
mixed human--agent--robot production.

\section{Synthesis of the Literature}

The literature provides five strong foundations:

\begin{enumerate}
\item automation changes the task allocation between capital and labour;
\item automation can affect wages, employment, inequality, and tax bases;
\item targeted taxation of automation may be optimal under some conditions;
\item social transfers and public capital ownership can distribute automation
gains; and
\item adaptive computational methods can support tax-policy design.
\end{enumerate}

Nevertheless, these literatures are generally connected only at a macroeconomic
or conceptual level. A practical fiscal architecture for an agentic economy
requires a missing intermediate layer: \textit{auditable attribution of
economic value to the productive factors that actually performed a task}.

\section{Research Gap}

Three unresolved gaps motivate CivicDividendOS.

\subsection{The Attribution Gap}

Existing robot-tax proposals frequently rely on robot counts, robot investment,
imputed wages, capital expenditure, or firm-level profits. These measures become
increasingly problematic in an agentic economy. One AI agent may execute the
work of many humans at negligible marginal software cost, while another AI
system may simply augment employees. A robot may perform a hazardous task that
should be encouraged rather than taxed. A cloud-based agent may serve multiple
jurisdictions simultaneously.

Therefore, neither ``one robot = one taxable worker'' nor a simple tax on
automation investment provides a reliable measure of the economic value created
by autonomous systems.

\subsection{The Liability Gap}

Giving AI agents or robots fictional salaries and asking them to pay income tax
creates legal and administrative difficulties. Machines do not independently
own their output under current systems. The economically relevant beneficiary
is generally an owner, firm, platform, or operator.

A new framework must therefore identify machine-originated economic activity
without making the machine itself the legal taxpayer.

\subsection{The Distribution Gap}

Current proposals usually choose one redistribution mechanism: a robot tax,
higher capital taxation, UBI, retraining, or public ownership. Yet an AI economy
can require several simultaneously. Displaced workers have immediate transition
needs; future workers need education; citizens may need a new non-wage source of
income; governments need stable revenue; and innovation still requires private
returns.

The literature does not yet provide an integrated, task-origin fiscal mechanism
that determines both \textit{how automation surplus enters the public system}
and \textit{how that surplus is divided among current citizens, affected
workers, future human capability, and public capital}.

\section{Problem Statement}

Let economic output $Y$ be jointly produced by human labour $H$, AI-agent
services $A$, robotic services $R$, data $D$, and conventional capital $K$:

\begin{equation}
Y = F(H,A,R,D,K).
\end{equation}

Under a traditional fiscal system, a substantial share of public revenue is
generated from:

\begin{equation}
Tax^{traditional}
=
\tau_L W_H + \tau_K \Pi + \tau_C C,
\end{equation}

where $W_H$ is human wage income, $\Pi$ represents taxable capital or corporate
income, and $C$ consumption.

As autonomous productive factors replace or augment human tasks, the wage share
may decline while output remains constant or increases:

\begin{equation}
\frac{W_H}{Y} \downarrow,
\qquad
\frac{Y_A+Y_R}{Y} \uparrow.
\end{equation}

If ownership of $A$ and $R$ is concentrated, household purchasing power,
payroll-financed social protection, and government revenue can weaken even when
aggregate productivity rises.

The fundamental research problem is therefore:

\begin{quote}
\textit{How can a tax-and-distribution system identify the economic value
created by humans, AI agents, robots, data, and conventional capital; assign
legal fiscal responsibility to the appropriate human or corporate
beneficiaries; and distribute part of automation-generated surplus in a way
that preserves innovation, public revenue, household income, workforce
transition, and broad participation in future capital returns?}
\end{quote}

The problem becomes more difficult when AI agents are cloud-based, operate
across borders, change tasks rapidly, and act primarily as complements rather
than substitutes for humans.

\section{Proposed Framework: CivicDividendOS}

CivicDividendOS contains seven interconnected components.

\subsection{1. Autonomous Economic Activity Passport}

Every economically significant autonomous deployment is assigned an
\textit{Autonomous Economic Activity Passport} (AEAP).

The passport is not legal personhood. It is a fiscal measurement record:

\begin{equation}
AEAP_j =
\{ID_j,O_j,J_j,T_j,E_j,Q_j,R_j,M_j\},
\end{equation}

where:
$O_j$ is the legally responsible owner/operator,
$J_j$ deployment jurisdiction,
$T_j$ task category,
$E_j$ energy/resource consumption,
$Q_j$ measured output,
$R_j$ risk classification, and
$M_j$ model/robot class.

The AEAP allows economic activity generated by an AI agent or robot to be
audited without pretending that the system itself is a taxpayer.

\subsection{2. Factor-Origin Ledger}

For every production episode $p$, CivicDividendOS estimates the marginal or
causal contribution of five productive factors:

\begin{equation}
\boldsymbol{\phi}_p =
\{\phi_H,\phi_A,\phi_R,\phi_D,\phi_K\},
\end{equation}

subject to

\begin{equation}
\sum_f \phi_f = 1.
\end{equation}

The attributable value generated by factor $f$ is:

\begin{equation}
V_{p,f}=\phi_{p,f} V_p.
\end{equation}

Contribution estimates may be obtained through controlled counterfactuals,
production-function estimation, Shapley-style contribution methods, matched
historical baselines, or digital-twin simulation.

This is a central departure from traditional robot taxation: the system taxes
neither the physical count of robots nor a fictional machine wage. It estimates
the economic origin of value.

\subsection{3. Substitution--Augmentation Classifier}

Automation should not be treated equally when it replaces a worker and when it
makes a worker more productive.

Define the automation mode:

\begin{equation}
Z_p \in
\{\text{Substitution},\text{Augmentation},
\text{NewTask},\text{SafetyReplacement}\}.
\end{equation}

A substitution score can be estimated as:

\begin{equation}
S_p =
\Pr(\Delta H_p < 0 \mid A_p,R_p,\mathbf{X}_p),
\end{equation}

while an augmentation score measures gains to human productivity or wages.

This classification feeds directly into the fiscal rate.

\subsection{4. Automation Contribution Base}

The proposed \textit{Automation Contribution Base} (ACB) is:

\begin{equation}
ACB_i =
VA^{A,R}_i - C^{allowable}_{A,R,i},
\end{equation}

where $VA^{A,R}_i$ is value added attributable to AI-agent and robotic factors
and $C^{allowable}$ represents verified operating, compute, energy,
depreciation, and acquisition costs according to policy.

The legal taxpayer is:

\begin{equation}
Taxpayer_i = Owner/Operator/Beneficiary(AEAP_i),
\end{equation}

not the AI agent or robot.

\subsection{5. Dynamic Social Automation Contribution}

The \textit{Social Automation Contribution} (SAC) is:

\begin{equation}
SAC_i = r_i \cdot ACB_i,
\end{equation}

where the adaptive rate $r_i$ is:

\begin{align}
r_i = r_0
&+\alpha S_i
+\beta C_i^{rent}
+\gamma E_i^{disp}
+\delta X_i^{ext}
+\eta R_i^{revenue} \nonumber\\
&-\theta A_i^{aug}
-\lambda T_i^{train}
-\mu B_i^{broad}
-\nu N_i^{newtask}.
\end{align}

Here:

\begin{itemize}
\item $S_i$ = labour substitution intensity;
\item $C_i^{rent}$ = economic-rent/concentration score;
\item $E_i^{disp}$ = measurable worker-displacement burden;
\item $X_i^{ext}$ = environmental or social externalities;
\item $R_i^{revenue}$ = erosion of payroll/social-contribution base;
\item $A_i^{aug}$ = human augmentation benefit;
\item $T_i^{train}$ = verified worker training and transition investment;
\item $B_i^{broad}$ = broad-based ownership/benefit score; and
\item $N_i^{newtask}$ = creation of new human tasks.
\end{itemize}

Thus automation that complements labour can face a low or even zero incremental
contribution, while high-rent labour-replacing automation can face a higher
rate.

\subsection{6. Human Earned-Income Shield}

A core innovation is to use automation revenue partly to reduce the tax wedge on
human work.

Let the labour-income tax relief rate be:

\begin{equation}
\Delta \tau_L(t) =
-\kappa
\frac{Rev_{SAC}(t)}{TaxBase_L(t)}.
\end{equation}

As autonomous production contributes more to public revenue, qualifying
low- and middle-income human workers receive lower payroll or earned-income
tax burdens.

This avoids a system in which society simultaneously taxes humans heavily for
working while giving machines a zero marginal social-contribution burden.

\subsection{7. Public Automation Wealth Fund}

A fixed fraction of SAC revenue and selected equity-based contributions enter a
professionally managed \textit{Public Automation Wealth Fund} (PAWF):

\begin{equation}
PAWF_{t+1}
=
PAWF_t(1+r_t^{fund})
+\omega Rev_{SAC,t}
+Equity_t
-Dist_t.
\end{equation}

The fund converts a portion of temporary AI and automation rents into long-lived
public capital.

Annual distributable returns are divided among four accounts:

\begin{equation}
Dist_t =
D^{citizen}_t+
D^{transition}_t+
D^{capability}_t+
D^{public}_t.
\end{equation}

\subsubsection{A. Civic Automation Dividend}

Every eligible resident receives a base dividend from investment returns rather
than directly consuming all annual tax receipts:

\begin{equation}
CADiv_t =
\frac{\rho D^{citizen}_t}{N_t}.
\end{equation}

This creates a non-wage household income stream linked to the productivity of
the automated economy.

\subsubsection{B. Worker Transition Insurance}

Workers experiencing verified displacement receive temporary wage insurance,
mobility support, or reskilling finance:

\begin{equation}
WTI_i =
g(WageLoss_i,Tenure_i,Exposure_i,Training_i).
\end{equation}

\subsubsection{C. Human Capability Endowment}

Part of the return finances education, apprenticeships, advanced technical
skills, public-service expertise, and human capabilities likely to complement
future AI.

\subsubsection{D. Public-Service Account}

A final portion supports social insurance and public services whose funding
previously depended heavily on payroll contributions.

\section{Automation Ownership Units}

CivicDividendOS adds a predistribution mechanism rather than relying exclusively
on ex-post tax transfers.

When a qualifying firm receives major public AI subsidies, public compute,
exclusive public data access, or government procurement support, a policy may
require a small \textit{Automation Ownership Unit} contribution:

\begin{equation}
AOU_i = \xi_i \cdot Equity_i.
\end{equation}

These units enter the PAWF rather than individual government departments.

This mechanism resembles the logic of a resource fund: the public contributes
infrastructure, research, data, procurement demand, education, and legal
institutions that help create private automation value, and therefore may hold
a modest diversified claim on future returns.

\section{Distribution Waterfall}

The complete income-distribution mechanism can be represented as:

\begin{equation}
Y = Y_H + Y_A + Y_R + Y_D + Y_K.
\end{equation}

Private primary distribution remains:

\begin{equation}
Y^{private} =
W_H + \Pi_{owners} + R_D + R_K.
\end{equation}

CivicDividendOS then introduces a secondary and ownership-based distribution:

\begin{align}
Y_H^{final} =&\;
W_H - Tax_H + Shield_H + Dividend_H \nonumber\\
&+Transition_H + Services_H,
\end{align}

while owners receive:

\begin{equation}
Y_O^{final}
=
\Pi_O - CIT_O - SAC_O + PrivateReturn_O.
\end{equation}

The state/public receives:

\begin{equation}
Y_P =
CIT + SAC + PAWFReturns,
\end{equation}

which is split among services, human-capability investment, transition support,
and citizen dividends.

The design therefore preserves private profit incentives while creating a
parallel public claim on part of the automation surplus.

\section{Fiscal Digital Twin}

Before changing SAC rates or distribution shares, policy is tested in a fiscal
digital twin containing heterogeneous households, firms, AI agents, robots,
sectors, and jurisdictions.

The policy vector is:

\begin{equation}
\Theta =
\{r_0,\alpha,\beta,\gamma,\omega,\rho,
\tau_L,\xi,\text{credits}\}.
\end{equation}

The policy engine seeks:

\begin{equation}
\Theta^{*}
=
\arg\max_{\Theta}
\left[
Growth + Revenue + Welfare + Employment
+ Innovation - Inequality - Leakage
\right]
\end{equation}

subject to fiscal-sustainability and minimum-innovation constraints.

The digital twin must report tax incidence explicitly. A nominal tax on an AI
operator may ultimately fall on shareholders, workers, suppliers, or consumers;
therefore the policy is not adjusted from statutory rates alone.

\section{Cross-Border Deployment Nexus}

Cloud AI complicates conventional tax residence because an agent can be trained
in one jurisdiction, owned in another, hosted in a third, and economically
deployed in a fourth.

CivicDividendOS therefore proposes a \textit{Deployment Nexus}:

\begin{equation}
Nexus_{ij}
=
w_1 Usage_{ij}
+w_2 Revenue_{ij}
+w_3 Users_{ij}
+w_4 PhysicalOps_{ij}.
\end{equation}

Automation contribution revenue can then be apportioned across jurisdictions
according to verified economic deployment rather than only server location.

This concept would require coordination with existing international corporate
tax rules and is proposed as a research mechanism rather than a claim of current
law.

\section{Illustrative Example}

Assume an urban logistics company previously employs 100 human dispatchers and
drivers. It introduces AI dispatch agents and autonomous delivery robots.

A conventional tax system observes:

\begin{itemize}
\item lower payroll;
\item fewer payroll contributions;
\item higher firm profit;
\item larger investment in automated capital.
\end{itemize}

A simple robot tax may charge each robot a fixed annual amount. This ignores
whether a robot replaced a worker, increased safety, or created new demand.

CivicDividendOS instead records task-origin activity.

Suppose the Factor-Origin Ledger estimates:

\[
\phi_H=0.30,\quad
\phi_A=0.25,\quad
\phi_R=0.30,\quad
\phi_D=0.05,\quad
\phi_K=0.10.
\]

The AI and robot share of attributable value is therefore:

\begin{equation}
VA^{A,R}=0.55V.
\end{equation}

The SAC rate is then reduced for verified safety improvements, creation of new
maintenance jobs, and retraining, but increased if labour displacement and
economic-rent concentration are high.

Revenue is not paid by the robots. It is paid by the firm that owns or
beneficially uses the automation.

Part of the revenue immediately funds transition insurance. Another portion
enters the PAWF. Investment returns then provide dividends to citizens and
finance future human skills.

The long-term allocation therefore becomes:

\[
\text{Human} \rightarrow
\text{Wages + Lower Labour Tax + Dividend},
\]

\[
\text{Owner} \rightarrow
\text{Profit - Corporate Tax - SAC + Investment Return},
\]

\[
\text{Public} \rightarrow
\text{SAC + Public Equity Returns},
\]

\[
\text{Future Workers} \rightarrow
\text{Human Capability Endowment}.
\]

\section{Research Questions}

\begin{itemize}

\item[\textbf{RQ1}] Can task-level factor-origin accounting estimate the value
attributable to humans, AI agents, robots, data, and conventional capital more
accurately than robot counts, imputed robot wages, or capital expenditure?

\item[\textbf{RQ2}] Can a substitution--augmentation classifier distinguish
automation that displaces taxable labour from automation that primarily
complements human workers?

\item[\textbf{RQ3}] Does a dynamic Social Automation Contribution preserve
government revenue with less distortion to innovation than a fixed robot tax?

\item[\textbf{RQ4}] Can combining automation contributions with a Human
Earned-Income Shield improve employment and after-tax wage outcomes relative to
raising labour or consumption taxes?

\item[\textbf{RQ5}] Does a Public Automation Wealth Fund create a more durable
and equitable distribution of AI-era wealth than transfers financed entirely
from annual tax revenue?

\item[\textbf{RQ6}] What allocation among citizen dividends, worker transition,
human-capability investment, and public services maximizes long-run social
welfare?

\item[\textbf{RQ7}] How should cross-border AI-agent activity be apportioned
when ownership, compute infrastructure, users, and economic deployment occur in
different jurisdictions?

\item[\textbf{RQ8}] How robust is CivicDividendOS to tax avoidance, model
misclassification, strategic task decomposition, automation offshoring, and
rapid changes in AI capability?

\end{itemize}

\section{Expected Contributions}

The proposed research aims to make the following contributions:

\begin{enumerate}

\item A \textit{Factor-Origin Ledger} for attributing value creation in mixed
human--AI--robot production.

\item An \textit{Autonomous Economic Activity Passport} that identifies machine
economic activity without granting machines fictional legal taxpayer status.

\item A \textit{Substitution--Augmentation Classifier} linking the fiscal
treatment of automation to its actual effect on human work.

\item An \textit{Automation Contribution Base} derived from attributable value
added rather than robot counts or fictional machine salaries.

\item A dynamic \textit{Social Automation Contribution} with explicit credits
for augmentation, training, new-task creation, and broad social benefit.

\item A \textit{Human Earned-Income Shield} that recycles part of automation
revenue into lower fiscal burdens on human work.

\item A \textit{Public Automation Wealth Fund} combining redistribution with
predistribution by building public ownership of productive AI-era capital.

\item A four-channel distribution waterfall comprising citizen dividends,
worker-transition insurance, human-capability investment, and public services.

\item A \textit{Deployment Nexus} for research into jurisdictional allocation of
cloud-based autonomous economic activity.

\item A fiscal digital twin for testing incidence, innovation, inequality,
employment, public revenue, and international leakage before policy deployment.

\end{enumerate}

\section{Evaluation Methodology}

The framework can be evaluated through a heterogeneous-agent macroeconomic and
urban-production simulation.

The simulated economy should contain:

\begin{itemize}
\item households with different income, skills, wealth, and automation exposure;
\item firms with heterogeneous productivity and market power;
\item conventional capital;
\item AI-agent capital;
\item robotic capital;
\item government tax and transfer accounts;
\item a public automation wealth fund; and
\item multiple domestic and cross-border sectors.
\end{itemize}

CivicDividendOS should be compared with:

\begin{enumerate}
\item existing labour/capital tax system;
\item fixed robot tax;
\item broad capital-income tax increase;
\item UBI funded through general taxation;
\item automation tax plus targeted retraining;
\item sovereign/public wealth fund without factor-origin taxation; and
\item proposed integrated CivicDividendOS.
\end{enumerate}

Evaluation metrics should include:

\begin{itemize}
\item GDP and productivity;
\item AI/robot adoption;
\item employment and hours;
\item wage share and capital share;
\item household disposable income;
\item Gini coefficients for income and wealth;
\item poverty rate;
\item government revenue volatility;
\item payroll/social-insurance funding;
\item investment and innovation;
\item firm entry and exit;
\item public-fund value;
\item citizen dividend per capita;
\item worker-transition duration;
\item human-capability investment;
\item tax avoidance/leakage; and
\item intergenerational wealth distribution.
\end{itemize}

Sensitivity analysis should vary AI costs, robot costs, substitution elasticities,
market concentration, international mobility of capital, tax compliance, and
the speed of technological progress.

\section{Policy Interpretation}

CivicDividendOS does not assume that robots or AI should always face additional
tax. It instead adopts four principles.

First, \textit{machines are measured but owners are taxed}. Second,
\textit{substitution and augmentation receive different fiscal treatment}.
Third, \textit{part of automation revenue is converted into broadly owned
capital rather than immediately spent}. Fourth, \textit{human work should not
remain the easiest tax base simply because it is administratively familiar}.

This produces a fiscal architecture that can evolve as the economy moves from
human-dominant production toward increasingly mixed or capital-dominant
production.

\section{Conclusion}

The AI and robotics era creates a distribution problem before it creates a
purely technical tax problem. If increasingly autonomous productive systems
generate a rising share of economic value while wages remain the principal
source of household income and payroll taxation remains central to public
finance, productivity can increase while purchasing power, social-insurance
finance, and wealth distribution become more fragile.

The proposed CivicDividendOS framework addresses this mismatch by separating
three questions that are frequently conflated: who created economic value, who
is legally responsible for tax, and who should ultimately share in the returns.

The Factor-Origin Ledger answers the first question. The Autonomous Economic
Activity Passport links the second to the legally responsible owner or operator.
The Social Automation Contribution, Human Earned-Income Shield, Public
Automation Wealth Fund, and four-channel distribution waterfall address the
third.

The proposed system therefore does not seek to make robots ``tax citizens.''
Instead, it creates a measurable and adaptive fiscal bridge between an economy
in which humans earn most productive income and one in which a larger share of
value may be generated by autonomous capital. Its long-term objective is to
ensure that technological productivity can remain privately innovative while
also becoming broadly investable, taxable, and socially shareable.

\begin{thebibliography}{99}

\bibitem{AcemogluRestrepo2019}
D. Acemoglu and P. Restrepo,
``Automation and new tasks: How technology displaces and reinstates labor,''
\textit{Journal of Economic Perspectives}, vol. 33, no. 2, pp. 3--30, 2019.

\bibitem{AcemogluRestrepo2020}
D. Acemoglu and P. Restrepo,
``Robots and jobs: Evidence from U.S. labor markets,''
\textit{Journal of Political Economy}, vol. 128, no. 6, pp. 2188--2244, 2020.

\bibitem{Dauth2021}
W. Dauth, S. Findeisen, J. S{\"u}dekum, and N. Woessner,
``The adjustment of labor markets to robots,''
\textit{Journal of the European Economic Association},
vol. 19, no. 6, pp. 3104--3153, 2021.

\bibitem{Zhang2019}
P. Zhang,
``Automation, wage inequality and implications of a robot tax,''
\textit{International Review of Economics \& Finance}, 2019.

\bibitem{Guerreiro2022}
J. Guerreiro, S. Rebelo, and P. Teles,
``Should robots be taxed?''
\textit{The Review of Economic Studies}, vol. 89, no. 1,
pp. 279--311, 2022.

\bibitem{Thuemmel2023}
U. Thuemmel,
``Optimal taxation of robots,''
\textit{Journal of the European Economic Association},
vol. 21, no. 3, pp. 1154--1190, 2023.

\bibitem{AcemogluOptimalPolicy}
D. Acemoglu, A. Manera, and P. Restrepo,
``Does the U.S. tax code favor automation?''
NBER Working Paper / related optimal-policy task-framework research, 2020.

\bibitem{InefficientAutomation}
``Inefficient automation,''
\textit{The Review of Economic Studies}, advance article,
doi: 10.1093/restud/rdae019.

\bibitem{Hotte2024}
K. H{\"o}tte, A. Theodorakopoulos, and P. Koutroumpis,
``Automation and taxation,''
\textit{Oxford Economic Papers}, vol. 76, no. 4,
pp. 945--973, 2024.

\bibitem{Merola2022}
R. Merola,
``Inclusive growth in the era of automation and AI:
How can taxation help?''
\textit{Frontiers in Artificial Intelligence}, vol. 5, 2022.

\bibitem{IMF2024AI}
International Monetary Fund,
``AI will transform the global economy. Let's make sure it benefits humanity,''
Jan. 2024.

\bibitem{IMF2024Fiscal}
International Monetary Fund,
``Fiscal policy can help broaden the gains of AI to humanity,''
June 2024.

\bibitem{IMF2026GlobalAI}
International Monetary Fund,
``Global economic and financial implications of artificial intelligence,''
2026.

\bibitem{OECD2025Korea}
OECD,
\textit{Artificial Intelligence and the Labour Market in Korea},
OECD Publishing, 2025.

\bibitem{Korinek2026}
A. Korinek et al.,
``Public finance in the age of AI: A primer,''
NBER Working Paper No. 34873, 2026.

\bibitem{FaivreCen2026}
J. Faivre and S. H. Cen,
``Taxing artificial intelligence,''
arXiv:2607.02144, 2026.

\bibitem{DimitropoulouAIIncome}
C. Dimitropoulou,
``Current income (profit) taxation of artificial intelligence,''
in research on robot and AI taxation, Edward Elgar Publishing.

\bibitem{ILOUBI}
International Labour Organization,
``Universal basic income proposals in light of ILO standards,''
ILO, Geneva.

\bibitem{LeGrand2020}
J. Le Grand,
``A springboard for new citizens: Universal basic capital,''
\textit{LSE Public Policy Review}, 2020.

\bibitem{Corneo2022}
G. Corneo,
``Progressive sovereign wealth funds,''
2022.

\bibitem{APF}
International Forum of Sovereign Wealth Funds,
``Alaska Permanent Fund,'' institutional profile.

\bibitem{BrookingsTax2026}
Brookings Institution,
``The future of tax policy: A public finance framework for the age of AI,''
Jan. 2026.

\bibitem{BrookingsFairness2026}
Brookings Institution,
``AI growth acceleration versus distributional fairness,''
May 2026.

\bibitem{Vincent2019}
N. Vincent, Y. Li, R. Zha, and B. Hecht,
``Mapping the potential and pitfalls of data dividends as a means of
sharing the profits of artificial intelligence,''
arXiv:1912.00757, 2019.

\bibitem{Bax2019}
E. Bax,
``Computing a data dividend,''
arXiv:1905.01805, 2019.

\bibitem{TaxAgent2025}
J. Wang, X. Fang, L. Huang, and Y. Huang,
``TaxAgent: How large language model designs fiscal policy,''
arXiv:2506.02838, 2025.

\bibitem{Fratric2025}
P. Fratri{\v c}, N. Holzenberger, and D. Restrepo Amariles,
``Can AI expose tax loopholes? Towards a new generation of legal policy
assistants,''
arXiv:2503.17339, 2025.

\end{thebibliography}

\end{document}
